\section{Условно математическо очакване}

\subsection{Мотивация и средноквадратично приближение}
В много практически ситуации разполагаме със случайна величина $X$, която описва дадена числова характеристика (например приходи, разходи на клиент, време на работа на процесор и т.н.). Вече знаем, че обикновеното математическо очакване $\E X$ е числото, което най-добре характеризира $X$ в средноквадратичен смисъл. Това означава, че дисперсията $\Var X = \E(X - \E X)^2$ е минималното средноквадратично отклонение от произволна константа $a$:
\[
\Var X = \min_{a \in \mathbb{R}} \E(X - a)^2
\]
Следователно, очакването е най-добрата константа за оценка на $X$, ако не разполагаме с никаква допълнителна информация.

Какво обаче се случва, ако разполагаме с допълнителна, наблюдаема случайна величина $Y$, която носи информация за $X$? Например, ако $X$ е разходът на клиент, а $Y$ е неговият пол (мъж или жена), е логично да очакваме, че оценката ни за $X$ може да се подобри, ако се съобразим с наблюдаваната стойност на $Y$. В такъв случай вече не търсим просто едно число $a$, което да минимизира средноквадратичната грешка, а функция $G(Y)$, зависеща от наблюдаваната стойност на $Y$. Целта ни е да намерим функция $G$, за която се достига следният минимум:
\[
\min_{G} \E(X - G(Y))^2
\]

Нека разгледаме прост модел, в който $Y$ е бинарна случайна величина, която разбива пространството от елементарни събития $\Omega$ на две части — събитието $A$ (например клиентът е мъж) и неговото допълнение $\overline{A}$ или $A^c$ (клиентът е жена). Можем да запишем $Y$ чрез индикаторната функция $\ind_A$:
\[
Y = \ind_A = \begin{cases}
1, & \text{ако се сбъдне } A \\
0, & \text{ако се сбъдне } A^c
\end{cases}
\]
Нека вероятностите са съответно $\mathbb{P}(Y=1) = \mathbb{P}(A) = p$ и $\mathbb{P}(Y=0) = \mathbb{P}(A^c) = 1-p = q$. Тъй като $Y$ приема само две стойности, всяка възможна функция $G(Y)$ от нея е изключително проста и може да бъде записана като линейна комбинация от индикаторите на $A$ и $A^c$ с неизвестни константи $a, b \in \mathbb{R}$:
\[
G(Y) = a \ind_A + b \ind_{A^c}
\]
Тогава оптимизационната ни задача се свежда до минимизиране на функцията на две променливи $f(a, b)$:
\[
\min_{a, b \in \mathbb{R}} \underbrace{\E(X - a \ind_A - b \ind_{A^c})^2}_{f(a,b)} = \min_{a, b \in \mathbb{R}} f(a,b)
\]
Разкриваме скобите, като използваме свойствата на индикаторите, че $\ind_A^2 = \ind_A$ и че тяхното произведение $\ind_A \ind_{A^c} = 0$:
\[
f(a,b) = \E[X^2 + a^2 \ind_A + b^2 \ind_{A^c} - 2aX\ind_A - 2bX\ind_{A^c}]
\]
Вземайки предвид, че $\E\ind_A = \mathbb{P}(A)$, диференцираме получената функция $f(a,b)$ спрямо променливите $a$ и $b$, и приравняваме частните производни на нула, за да намерим минимума:
\[
0 = \frac{\partial f}{\partial a} = 2a\mathbb{P}(A) - 2\E[X\ind_A] \implies a = \frac{\E[X\ind_A]}{\mathbb{P}(A)} = \frac{\E[XY]}{\mathbb{P}(Y=1)}
\]
\[
0 = \frac{\partial f}{\partial b} = 2b\mathbb{P}(A^c) - 2\E[X\ind_{A^c}] \implies b = \frac{\E[X\ind_{A^c}]}{\mathbb{P}(A^c)} = \frac{\E[X(1-Y)]}{\mathbb{P}(Y=0)}
\]
Полученият резултат е минимум (тъй като изразът представлява квадратична форма с положителен хесиан). Функцията, за която се достига минимумът, отбелязваме с $G^*(Y)$:
\[
G^*(Y) = \frac{\E[X\ind_A]}{\mathbb{P}(A)} \ind_A + \frac{\E[X\ind_{A^c}]}{\mathbb{P}(A^c)} \ind_{A^c} =: \E(X\given Y)
\]
Тази функция дефинира \emph{условното математическо очакване} на $X$ при условие $Y$. Вместо да използваме едно единствено число, вече използваме две различни стойности в зависимост от това каква стойност сме наблюдавали за $Y$. Ако клиентът е мъж, използваме оценката $a$, ако е жена — оценката $b$.

За да илюстрираме как условната вероятност естествено възниква в този контекст, нека приемем, че $X$ също е бинарна случайна величина, например $X = \ind_B$. Тогава $\E[X\ind_A] = \E[\ind_B \ind_A] = \mathbb{P}(A \cap B)$. Замествайки във формулата, получаваме:
\[
G^*(Y) = \frac{\mathbb{P}(A \cap B)}{\mathbb{P}(A)} \ind_A + \frac{\mathbb{P}(A^c \cap B)}{\mathbb{P}(A^c)} \ind_{A^c} = \mathbb{P}(B\given A)\ind_A + \mathbb{P}(B\given A^c)\ind_{A^c}
\]
Виждаме, че коефициентите пред индикаторите са точно условните вероятности $\mathbb{P}(B\given A)$ и $\mathbb{P}(B\given A^c)$.

\subsection{Дефиниция за дискретни случайни величини}

Въз основа на изложената геометрична мотивация (като проекция в средноквадратичен смисъл), можем да дефинираме понятието формално за произволни дискретни случайни величини.

\begin{keydefn}
Нека $X$ и $Y$ са дискретни случайни величини. Условното математическо очакване на $X$ при условие, че $Y$ е приело конкретната стойност $y_j$, се дефинира като сума по всички възможни стойности на $X$ с техните условни вероятности:
\[
\E[X \given Y=y_j] = \sum_i x_i \mathbb{P}(B_i \given A_j)
\]
където $B_i = \{X=x_i\}$ са възможните стойности на $X$, а $A_j = \{Y=y_j\}$.
\end{keydefn}

Това е просто число (за дадено фиксирано събитие $Y=y_j$). Ако обаче разглеждаме $Y$ като случайна величина, можем да дефинираме \emph{условното математическо очакване на $X$ спрямо $Y$} (което бележим с $\E[X\given Y]$) като нова дискретна случайна величина, чиято стойност зависи изцяло от това каква стойност ще приеме $Y$:
\[
\E[X \given Y] = \sum_j \E[X \given Y=y_j] \ind_{A_j}
\]

\begin{prop}
Ако $X$ и $Y$ са дискретни случайни величини, то $\E[X\given Y]$ е също дискретна случайна величина, която има следната таблица на разпределение:
\[
\begin{array}{c|c|c|c|c}
\E[X \given Y] & \E[X \given Y=y_1] & \dots & \E[X \given Y=y_j] & \dots \\
\hline
\mathbb{P} & p_1 = \mathbb{P}(A_1) & \dots & p_j = \mathbb{P}(A_j) & \dots
\end{array}
\]
Тази случайна величина изцяло се предопределя от стойностите на $Y$ и приема съответната условна очаквана стойност със същата вероятност, с която $Y$ приема стойността $y_j$.
\end{prop}

\subsection{Свойства на условното математическо очакване}

Условното очакване притежава серия от важни свойства, аналогични на обикновеното математическо очакване. Те са валидни както за дискретни, така и за непрекъснати величини, но тук ще бъдат обосновани интуитивно през призмата на дискретните величини.

\begin{prop}[Свойства на условното математическо очакване]\label{prop:cond-exp-props}
\begin{enumerate}
    \item[а)] \textbf{Линейност:}
    $\E[aX + bZ \given Y] = a\E[X \given Y] + b\E[Z \given Y]$.
    \item[б)] \textbf{Независимост:} ако $X \perp\!\!\!\perp Y$, то $\E[X \given Y] = \E X$.
    \item[в)] \textbf{Детерминираност:} ако $X = f(Y)$, то $\E[X \given Y] = X$.
    \item[г)] \textbf{Повторно очакване (Закон за пълното математическо очакване):}
    $\E[\E[X \given Y]] = \E X$.
    \item[д)] \textbf{Изнасяне на функция от условието:}
    $\E[f(X,Y) \given Y=y_j] = \E[f(X, y_j) \given Y=y_j]$.
\end{enumerate}
\end{prop}

\begin{proof}
а) Използвайки дефиницията за индикаторно представяне:
\begin{align*}
\E[aX + bZ \given Y]
  &= \sum_j \frac{\E[(aX+bZ)\ind_{A_j}]}{\mathbb{P}(A_j)}\ind_{A_j} \\
  &= a \sum_j \frac{\E[X\ind_{A_j}]}{\mathbb{P}(A_j)}\ind_{A_j}
   + b \sum_j \frac{\E[Z\ind_{A_j}]}{\mathbb{P}(A_j)}\ind_{A_j} ,
\end{align*}
което е точно $a\E[X \given Y] + b\E[Z \given Y]$.

б) Тъй като $X$ и $\ind_{A_j}$ са независими, то $\E[X\ind_{A_j}] = \E X \cdot \E\ind_{A_j} = \E X \cdot \mathbb{P}(A_j)$. Тогава в дефиниционната сума вероятностите $\mathbb{P}(A_j)$ се съкращават и остава $\E X \sum_j \ind_{A_j} = \E X \cdot 1 = \E X$ (защото индикаторите на пълна група от събития се сумират до 1).

г) Прилагаме оператора $\E$ върху дефиниционната сума на $\E[X\given Y]$:
\begin{align*}
\E[\E[X\given Y]]
  &= \E\left[\sum_j \frac{\E[X\ind_{A_j}]}{\mathbb{P}(A_j)} \ind_{A_j}\right]
   = \sum_j \frac{\E[X\ind_{A_j}]}{\mathbb{P}(A_j)} \E\ind_{A_j} \\
  &= \sum_j \E[X\ind_{A_j}]
   = \E\left[\sum_j X\ind_{A_j}\right] = \E X
\end{align*}
\end{proof}

Свойства в) и д) са непосредствени. За в): ако познавайки $Y$ ние вече знаем с абсолютна точност стойността на $X$, то случайността отпада и най-добрата оценка е самото $X$. За д): тъй като в условието знаем, че $Y=y_j$, можем да заместим $Y$ с тази константа в самата функция.

Свойство б) има ясен смисъл: когато $Y$ не носи информация за $X$, условното очакване съвпада с безусловното, защото най-доброто число за приближение на $X$ (в средноквадратичен смисъл) остава просто нейното обикновено математическо очакване.

\begin{example}[Смес от разпределения]\label{ex:08-1}

Условното математическо очакване е изключително полезно при т.нар. смеси (mixture distributions), където параметрите на дадено разпределение самите те са случайни величини. Нека разгледаме Бернулиева случайна величина $X \sim \Ber(Y)$, чийто параметър на успеха е случайната величина $Y$:
\[
Y = \begin{cases}
\frac{2}{3}, & \text{с вероятност } p \\
\frac{1}{3}, & \text{с вероятност } 1-p
\end{cases}
\]
Това означава, че първо теглим параметъра $Y$. Ако се падне $2/3$, тогава $X$ е Бернулиева с вероятност за успех $2/3$. Ако $Y$ е $1/3$, тогава $X$ е Бернулиева с вероятност за успех $1/3$. Търсим пълната вероятност $\mathbb{P}(X=1)$, която за индикаторна променлива напълно съвпада с нейното математическо очакване $\E X$. Използвайки свойството за повторно очакване $\E X = \E[\E[X \given Y]]$, имаме:
\begin{align*}
\E[X \given Y] &= \E\left[X \given Y=\frac{2}{3}\right] \ind_{\{Y=2/3\}} + \E\left[X \given Y=\frac{1}{3}\right] \ind_{\{Y=1/3\}} \\
&= \frac{2}{3} \ind_{\{Y=2/3\}} + \frac{1}{3} \ind_{\{Y=1/3\}}
\end{align*}
Вземайки безусловното математическо очакване от двете страни, получаваме крайната вероятност за успех:
\[
\mathbb{P}(X=1) = \E[\E[X \given Y]] = \frac{2}{3} \mathbb{P}\left(Y=\frac{2}{3}\right) + \frac{1}{3} \mathbb{P}\left(Y=\frac{1}{3}\right) = \frac{2}{3}p + \frac{1}{3}(1-p)
\]

\end{example}

\section{Условни разпределения}

Аналогично на условното очакване, можем да дефинираме цялостното \emph{условно разпределение} на една случайна величина спрямо друга.

\begin{defn}
Нека $X$ и $Y$ са дискретни случайни величини. Условното разпределение на $X$ при условие $Y=y_j$ (за всяка възможна стойност на $Y$ с $\mathbb{P}(Y=y_j)>0$) се задава чрез условните вероятности за всяка възможна стойност $x_i$:
\[
\begin{array}{c|c|c|c|c|c}
X \given Y=y_j & x_1 & x_2 & \dots & x_i & \dots \\
\hline
\mathbb{P} & \mathbb{P}(X=x_1 \given Y=y_j) & \dots & \dots & \mathbb{P}(X=x_i \given Y=y_j) & \dots
\end{array}
\]
Това на практика са изменените вероятности за $X$, които произтичат от допълнителната информация, че е настъпило събитието $Y=y_j$.
\end{defn}

\begin{example}\label{ex:08-2}
Да разгледаме хвърляне на два зара. Нека $Y$ е броят на падналите се шестици, а $X$ е броят на падналите се единици. И двете случайни величини приемат стойности в множеството $\{0, 1, 2\}$. Съвместното им разпределение е зададено със следната таблица (известна от предходна лекция):
\[
\begin{array}{c|c|c|c}
Y \setminus X & 0 & 1 & 2 \\
\hline
0 & 16/36 & 8/36 & 1/36 \\
1 & 8/36 & 2/36 & 0 \\
2 & 1/36 & 0 & 0
\end{array}
\]
Искаме да намерим условното разпределение на $X$ при условие, че $Y=0$. Маргиналната вероятност за $Y=0$ е $\mathbb{P}(Y=0) = \frac{16}{36} + \frac{8}{36} + \frac{1}{36} = \frac{25}{36}$. За да намерим условните вероятности $\mathbb{P}(X=x_i \given Y=0)$, просто разделяме съвместната вероятност от таблицата на маргиналната вероятност:
\begin{align*}
\mathbb{P}(X=0 \given Y=0) &= \frac{16/36}{25/36} = \frac{16}{25} \\
\mathbb{P}(X=1 \given Y=0) &= \frac{8/36}{25/36} = \frac{8}{25} \\
\mathbb{P}(X=2 \given Y=0) &= \frac{1/36}{25/36} = \frac{1}{25}
\end{align*}
Този резултат има напълно ясна интуиция. Ако знаем със сигурност, че $Y=0$ (не се е паднала нито една шестица), това е еквивалентно на това да хвърляме зарове само с по 5 стени (защото сме елиминирали напълно шестицата). Възможните изходи от хвърлянето на два зара намаляват до $5^2 = 25$. Вероятността да не се падне единица е $\left(\frac{4}{5}\right)^2 = \frac{16}{25}$, да се падне точно една единица е $2 \cdot \frac{1}{5} \cdot \frac{4}{5} = \frac{8}{25}$, а две единици - $\left(\frac{1}{5}\right)^2 = \frac{1}{25}$.

\end{example}

\section{Полиномно разпределение}

Полиномното разпределение е естествено многомерно обобщение на биномното разпределение. Докато при биномното разпределение имаме само два възможни изхода (успех и неуспех), при полиномното разполагаме с $r \ge 1$ възможни изхода.
Нека провеждаме $n$ независими експеримента. Всеки експеримент може да завърши с един от $r$ изхода. Съответните вероятности за тези изходи са $p_1, p_2, \dots, p_r$, като сумата им естествено е $\sum_{j=1}^{r} p_j = 1$.

Интересува ни съвместното разпределение на вектора от случайни величини $(X_1, X_2, \dots, X_r)$, където $X_j$ е броят на настъпванията на изход $j$. Зададени са цели неотрицателни числа $k_1, k_2, \dots, k_r$, такива че сборът им дава общия брой експерименти: $k_1 + k_2 + \dots + k_r = n$.
Вероятността да се получи точно тази конфигурация от резултати се пресмята по следната формула:
\[
\mathbb{P}(X_1=k_1, X_2=k_2, \dots, X_r=k_r) = \frac{n!}{k_1! \dots k_r!} p_1^{k_1} p_2^{k_2} \dots p_r^{k_r}
\]
Тази формула се извежда аналогично на биномното разпределение. Вероятността за една конкретна наредба на тези изходи (т.е. конкретна последователност от $n$ събития) е произведението $p_1^{k_1} \dots p_r^{k_r}$. Тъй като нас не ни интересува конкретният ред на настъпване на събитията, а само общият им брой, трябва да умножим по броя на всички възможни конфигурации, които водят до този краен резултат.
Броят на начините да изберем $k_1$ места за първия изход е $\binom{n}{k_1}$. За останалите $n-k_1$ свободни места избираме $k_2$ места за втория изход, което може да стане по $\binom{n-k_1}{k_2}$ начина, и така нататък. Произведението от тези биномни коефициенти се опростява алгебрично до мултиномния коефициент:
\[
\binom{n}{k_1} \binom{n-k_1}{k_2} \dots \binom{n-k_1-\dots-k_{r-1}}{k_r} = \frac{n!}{k_1! \dots k_r!}
\]
Ако $r=2$, получаваме класическото биномно разпределение, където $k_2 = n - k_1$ и $p_2 = 1 - p_1$.

\section{Непрекъснати случайни величини}

\subsection{Въведение и мотивация}
Досега разглеждахме дискретни случайни величини, които могат да приемат най-много изброимо множество от стойности. Има обаче много случайни характеристики от реалния живот (например температура, тегло, време на живот на процесор, приходи), които чисто теоретично могат да приемат неизброимо много стойности — всяка стойност от някакъв реален интервал.

От гледна точка на компютърното моделиране, неизброимостта е фикция — ние винаги съхраняваме данните с крайна точност (чрез тактовете на процесора). В такъв случай защо въобще са ни нужни непрекъснати случайни величини?
Отговорът се крие в ефективността. Ако искаме да моделираме равномерна случайна величина в интервала $(0,1)$ с много голяма точност (например 100 милиона възможни стойности), трябва да пазим масив със 100 милиона вероятности, въпреки че всички те са напълно еднакви. Много по-изчислително ефективно е да използваме една проста функция (плътност), която компактно да описва поведението на модела.
Освен това, на практика изключително рядко се интересуваме от точната вероятност температурата да бъде съвършено равняваща се на дадено число с безкрайна точност (която вероятност клони към нула). Много по-често ни интересува вероятността стойността да попадне в даден диапазон.

\subsection{Дефиниция и плътност на разпределение}

\begin{defn}[Абсолютно непрекъсната случайна величина]
Случайната величина $X$ се нарича непрекъсната, ако съществува неотрицателна функция $f_X: \mathbb{R} \to [0, \infty)$, наричана \emph{плътност на разпределение} (probability density function), такава че:
\begin{enumerate}
    \item $f_X(x) \ge 0 \quad \forall x \in \mathbb{R}$
    \item Лицето под цялата крива е единица: $\int_{-\infty}^{\infty} f_X(x) dx = 1$
    \item За всеки две числа $a < b$ (включително $\pm\infty$) е изпълнено:
    \[
    \mathbb{P}(X \in (a, b)) = \int_{a}^{b} f_X(x) dx
    \]
\end{enumerate}
\end{defn}

Графиката на функцията $f_X(x)$ ни показва къде случайните величини са по-склонни да приемат стойности. Лицето под графиката между точките $a$ и $b$ е точно равно на вероятността $X$ да попадне в интервала $(a,b)$.

От свойствата на интеграла веднага следва, че вероятността $X$ да приеме точно една конкретна стойност $c$ е нула:
\[
\mathbb{P}(X = c) = \int_{c}^{c} f_X(x) dx = 0
\]
Поради това за непрекъснати случайни величини границите на интервала нямат никакво значение при пресмятане на вероятност:
\[
\mathbb{P}(a < X < b) = \mathbb{P}(a \le X < b) = \mathbb{P}(a < X \le b) = \mathbb{P}(a \le X \le b)
\]

\begin{example}[Медицински застраховки]\label{ex:08-3}
Нека разходите за медицински застраховки в една фирма са зададени с модела $M = 100\,000 \cdot X$, където $X$ е непрекъсната случайна величина с плътност:
\[
f_X(x) = \begin{cases}
5(1-x)^4, & x \in (0,1) \\
0, & \text{иначе}
\end{cases}
\]
Тази плътност е нула за отрицателни стойности и за стойности по-големи от 1. В интервала $(0,1)$ представлява намаляваща крива (започваща от стойност 5 при $x=0$). Търсим вероятността разходите на фирмата да надвишат 10 000 единици:
\[
\mathbb{P}(M > 10\,000) = \mathbb{P}\left(100\,000 \cdot X > 10\,000\right) = \mathbb{P}\left(X > \frac{1}{10}\right)
\]
Тъй като $X$ приема стойности само до 1, това се свежда до определен интеграл:
\[
\mathbb{P}\left(\frac{1}{10} < X < 1\right) = \int_{1/10}^{1} 5(1-x)^4 dx
\]
Този интеграл се решава изключително лесно и е равен на $\left(\frac{9}{10}\right)^5 \approx 0{,}59$.
Вместо да пази огромен масив с дискретни вероятности за всяка възможна стотинка, фирмата използва компактна функция за оценка на риска.

\end{example}

\subsection{Функция на разпределение}

Функцията на разпределение $F_X(x)$ за непрекъсната случайна величина $X$ се дефинира по обичайния начин като вероятността $X$ да бъде по-малко от $x$:
\[
F_X(x) = \mathbb{P}(X < x) = \int_{-\infty}^{x} f_X(y) dy \quad \forall x \in \mathbb{R}
\]
Това представлява натрупване (акумулиране) на площта под плътността от минус безкрайност до текущата точка $x$. Диференциалът $dy$ в този контекст трябва да се разбира като мярка.
Ако плътността $f_X$ е непрекъсната функция в дадена точка $x_0$, то от основната теорема на интегралното смятане следва, че производната на функцията на разпределение е равна на плътността:
\[
\left. \frac{dF_X}{dx} \right|_{x=x_0} = f_X(x_0) \implies F_X' = f_X
\]
Следователно, ако знаем плътността, чрез интегриране възстановяваме функцията на разпределение, а ако знаем функцията на разпределение, чрез диференциране възстановяваме плътността.

\section{Смяна на променливите при непрекъснати случайни величини}

Често имаме даден модел със случайна величина $X$, но се интересуваме от някаква нейна трансформация $Y = g(X)$ (например ако $X$ е приходът на клиент, $Y$ може да бъде неговият разход за конкретни мероприятия, който зависи функционално от прихода). Искаме да намерим плътността на новата случайна величина $Y$.

\begin{thm}
Нека $X$ е непрекъсната случайна величина, а $g: \mathbb{R} \to \mathbb{R}$ е строго монотонно растяща или строго монотонно намаляваща функция (поне в областта, където $X$ приема стойности с положителна плътност). Тогава $Y = g(X)$ е също непрекъсната случайна величина, чиято плътност се задава от формулата:
\[
f_Y(y) = f_X(g^{-1}(y)) \left| (g^{-1}(y))' \right|
\]
\end{thm}

\begin{proof}
Търсим функция $f_Y$, която при интегриране над произволен интервал $(a, b)$ да дава вероятността $\mathbb{P}(Y \in (a, b))$.
Разглеждаме първо случая, когато $g$ е строго монотонно растяща ($g^{\uparrow}$).
\[
\mathbb{P}(Y \in (a, b)) = \mathbb{P}(g(X) \in (a, b)) = \mathbb{P}(X \in (g^{-1}(a), g^{-1}(b)))
\]
Тъй като познаваме плътността на $X$, тази вероятност е просто интеграл:
\[
\mathbb{P}(Y \in (a, b)) = \int_{g^{-1}(a)}^{g^{-1}(b)} f_X(w) dw
\]
За да превърнем границите на интегриране отново в $a$ и $b$, правим субституцията $w = g^{-1}(v)$. Тогава диференциалът е $dw = (g^{-1}(v))' dv$. Границите на интегриране се променят от $g^{-1}(a)$ до $g^{-1}(b)$ съответно на $a$ до $b$:
\[
\int_{a}^{b} f_X(g^{-1}(v)) (g^{-1}(v))' dv
\]
Тъй като $g$ е монотонно растяща функция, нейната обратна функция $g^{-1}$ също е монотонно растяща, което означава, че производната $(g^{-1}(v))'$ е положителна, т.е. съвпада със своя модул.

Нека сега разгледаме случая, когато $g$ е строго монотонно намаляваща ($g^{\downarrow}$). Когато приложим $g^{-1}$ към неравенството $a < g(X) < b$, посоките на неравенствата се обръщат, защото функцията обръща наредбата:
\[
\mathbb{P}(g(X) \in (a, b)) = \mathbb{P}(X \in (g^{-1}(b), g^{-1}(a))) = \int_{g^{-1}(b)}^{g^{-1}(a)} f_X(w) dw
\]
При същата субституция $w = g^{-1}(v)$ границите на интегриране стават от $b$ до $a$. За да обърнем нормално интеграла да бъде от $a$ до $b$, трябва да добавим знак минус:
\[
-\int_{a}^{b} f_X(g^{-1}(v)) (g^{-1}(v))' dv
\]
Но тъй като $g$ е монотонно намаляваща, $(g^{-1}(v))'$ е отрицателна функция, следователно $-(g^{-1}(v))' = \left| (g^{-1}(v))' \right|$. И в двата случая стигаме до един и същ обединен израз:
\[
\int_{a}^{b} f_X(g^{-1}(v)) \left| (g^{-1}(v))' \right| dv
\]
Тъй като вероятността $\mathbb{P}(Y \in (a, b))$ се пресмята чрез интеграл над интервала $(a,b)$ от тази конкретна функция, то по дефиниция подинтегралната функция е именно търсената плътност на $Y$.
\end{proof}

\section{Числени характеристики на непрекъснатите величини}

\subsection{Математическо очакване и дисперсия}

По аналогичен начин на дискретните случайни величини, където математическото очакване се получаваше като сума от произведенията на стойностите с техните вероятности $\E X = \sum x_i p_i$, при непрекъснатите величини сумата логично се заменя с интеграл (тъй като интегралът е граница на сума), а вероятностите преминават в плътност.

\begin{defn}
Нека $X$ е непрекъсната случайна величина. Ако интегралът $\int_{-\infty}^{\infty} |x| f_X(x) dx < \infty$ е краен, то \emph{математическото очакване} на $X$ се дефинира като:
\[
\E X = \int_{-\infty}^{\infty} x f_X(x) dx
\]
От гледна точка на механиката, ако разглеждаме плътността $f_X(x)$ като разпределение на масата върху пръчка, очакването $\E X$ е точно центърът на тежестта на тази пръчка.
\end{defn}

\begin{defn}[дисперсия на непрекъсната случайна величина]\label{def:var-continuous}
\emph{Дисперсията} $\Var X$ продължава да се дефинира като средноквадратичното отклонение от очакването:
\[
\Var X = \E(X - \E X)^2 = \int_{-\infty}^{\infty} (x - \E X)^2 f_X(x) \, dx .
\]
\end{defn}

Ако искаме да пресметнем математическото очакване на функция от случайната величина $g(X)$, не е необходимо първо да намираме плътността на $Y=g(X)$. Използваме директната формула (която се доказва строго чрез теория на мярката, но следва интуитивно от току-що доказаната теорема за смяна на променливите):
\[
\E[g(X)] = \int_{-\infty}^{\infty} g(x) f_X(x) dx
\]

\textbf{Свойства на очакването и дисперсията:}
Следните свойства остават изцяло в сила както за дискретни, така и за непрекъснати случайни величини (нека $a, b, c$ са константи):
\begin{itemize}
    \item Очакване на константа: $\E c = c$
    \item Дисперсия на константа: $\Var c = 0$
    \item Линейност на очакването: $\E(aX + b) = a\E X + b$ и $\E(aX + bY) = a\E X + b\E Y$
    \item Изместване на дисперсията: $\Var(X + c) = \Var X$
    \item Мащабиране на дисперсията: $\Var(aX) = a^2 \Var X$
    \item Независимост: Ако $X \perp\!\!\!\perp Y$, то $\E[XY] = \E X \cdot \E Y$ и $\Var(X + Y) = \Var X + \Var Y$
\end{itemize}

\subsection{Равномерно разпределение}

Равномерното разпределение (Uniform distribution) е най-простият клас непрекъснати случайни величини. То се използва изключително често като априорно разпределение, когато нямаме никаква допълнителна информация или експертни познания за системата и искаме да избегнем субективизъм (т.е. стартираме на чисто).

Казваме, че $X$ е равномерно разпределена в интервала $[a, b]$, където $a < b$ са крайни числа, ако плътността ѝ е постоянна в този интервал и нула извън него. За да бъде лицето под графиката на плътността равно на 1 (както изисква дефиницията), плътността трябва да е:
\[
f_X(x) = \begin{cases}
\frac{1}{b-a}, & x \in [a, b] \\
0, & \text{иначе}
\end{cases}
\]
Лесно се проверява, че $\int_{a}^{b} \frac{1}{b-a} dx = \frac{1}{b-a} \cdot (b-a) = 1$.
Нарича се равномерно, защото за подинтервали с еднаква дължина вероятността случайната величина да попадне в тях е напълно еднаква, без значение къде се намира подинтервалът в рамките на $[a, b]$.

Функцията на разпределение за равномерното разпределение се получава чрез директно интегриране:
\[
F_X(x) = \begin{cases}
0, & x \le a \\
\frac{x-a}{b-a}, & x \in (a, b) \\
1, & x \ge b
\end{cases}
\]
Ако $x \in (a, b)$, то $F_X(x) = \int_{-\infty}^{x} f_X(y) dy = \int_{a}^{x} \frac{1}{b-a} dy = \frac{x-a}{b-a}$.
За $x \le a$ интегралът е от 0, следователно $0$. За $x \ge b$ плътността след $b$ отново е 0, така че целият интеграл се свежда до площта над целия интервал $[a, b]$, която е 1.

\section{Задачи}

\begin{enumerate}
    \item Нека $Y=aX+b$, където $a\ne0$. Използвайте теоремата за смяна на променливите, за да намерите плътността на $Y$ чрез плътността на $X$.
    \item Проверете от дефиницията, че функцията
    \[
    f(x)=\frac{1}{b-a}\ind_{[a,b]}(x),\qquad a<b,
    \]
    е плътност: тя е неотрицателна и интегралът ѝ върху реалната права е единица.
\end{enumerate}
